\subsection{Defining Rational Exponents: p/q} 
We have now defined the following exponents:
\begin{definition}
For positive integers $n$ and any real number $x$,
$$x^n=\underbrace{x\cdot x\cdots x}_{n\ \mathrm{copies}}$$
For any real number $x\neq 0$,
$$x^0=1$$
For any negative integer $-n$ and real number $x\neq 0$,
$$x^{-n}=\dfrac{1}{x^n}$$
For any positive integer $q$ and any real number $x$ for which $\sqrt[q]{x}$ is defined,
$$x^{1/q}=\sqrt[q]{x}$$
\end{definition}
We define rational exponents $\frac{p}{q}$ to be consistent with the power of a power rule.
\begin{lpic}{}{5}
\foreach \y/\lbl in {1/{Note that $\frac{p}{q}={}$},3/{$x^{p/q}={}$},5/{$x^{p/q}={}$}}
	\regtextleft{0.25}{-\y}{\lbl};
\end{lpic}
\begin{definition}
For any integers $p$ and $q$ with $q>0$, and any real number $x$ for which $\sqrt[q]{x}$ is defined,
$$x^{p/q}=\sqrt[q]{x^p}=\left(\sqrt[q]{x}\right)^p$$
\end{definition}